% page 473 of the facsimile (image p-472 of the scan)
% block 162.6 x 259.3 mm ; left margin 39.4 mm
\pgc{17.780mm}{0.000mm}{
\l{}
\l{}
\l{}
\l{\cel{47}465}
\l{}
\l{\sou{pronacar}. (\sou{trans.}) (\sou{anat.}) Igar la karpo en posturo tala ke}
\l{\cel{3}la palmo di la manuo \textquotedbl{}regardas\textquotedbl{} vers sulo. (antonimo :}
\l{\cel{3}\sou{supinacar}).-}
\l{}
\l{\sou{pronomo.}\cel{1}(\sou{gram.)}\cel{2}Vorto-speco qua remplasas e pleas la rolo di}
\l{\cel{3}la nomo tacita. DEFIS.}
\l{}
\l{\sou{pronta.}\cel{2}I. Tote preparita. - II. Qua esas quik....onta. - IS.}
\l{}
\l{\sou{pronukleo}. (\sou{biol.})\cel{2}Parto nuklea di la ovulo e di la spermato-}
\l{\cel{3}zoido qui kunfuzesas lor la fekundigo. -}
\l{}
\l{\sou{pronuncar.}\cel{2}(\sou{trans.}) I. Igar audebla publike (pronuncar diskur-}
\l{\cel{3}so, laudo funerala, pledo). - II. Enuncar, artikulante la}
\l{\cel{3}silabi, la vorti. - EFIS.}
\l{}
\l{\sou{propagar}. (\sou{trans.}) Diskonocigar proximope. - (\sou{metaf}.: genitar}
\l{\cel{3}por plumultigar (sua) decendanti.) - DEFIRS.}
\l{}
\l{proporcionar. (\sou{trans.}) Konformigar a la legi segun qui(lu),}
\l{\cel{3}kom parto di ensemblo, esas tale o tale granda, konsidere di}
\l{\cel{3}la toto. - DEFIRS.}
\l{}
\l{\sou{proporciono}. I. En linguo vulgara : (a) grandeso di parto}
\l{\cel{3}kompare a la toto od a la altra parti, en snsemblo ; (b)}
\l{\cel{3}grandeso di kozo kompare a kozo analoga, egardata kom tipo ;}
\l{\cel{3}(c) grandeso di kozo kompare ad altra kozo, la unesma kres-}
\l{\cel{3}kante o deskreskante segun ke la duesma kreskas o deskreskas.}
\l{\cel{3}- II. En la linguo ciencala : (a) (\sou{matem.)Proporciono geome}-}
\l{\cel{3}\sou{triala}\cel{1}: Inter-egaleso di du raporti per quociento.\cel{2}\sou{Propor}-}
\l{\cel{3}\sou{ciono-regulo}\cel{1}: operaco per qua, donite un ek la du raporti}
\l{\cel{3}ed un termino di la duesma, onu determinas la altra termino.}
\l{\cel{3}\sou{Prpporciono aritmetikala}\cel{1}: Interegaleso di du raporti per}
\l{\cel{3}difero. - (\sou{kemio})\cel{1}\sou{Proporciono definita}\cel{1}: Quantesi fixa, na-}
\l{\cel{3}variiva, segun qui interunionas la korpi por formacar kom-}
\l{\cel{3}binuri kemiala. - DEFIRS.}
\l{}
\l{\sou{propoza}r (\sou{trans}.)(ad). I. Pozar (ulo) avan ulu kom objekto, kom}
\l{\cel{3}temo, di studio, di exploro. - II. Pozar (ulo) avan ulu kom}
\l{\cel{3}regulo, kom modelo a qui konformigor sua agado. - III. Pozar}
\l{\cel{3}(ulu) avan (ulu, korporaciono) kom kandidato. - DEFIRS.}
\l{}
\l{\sou{propoziciono}. I. (\sou{logiko}) Enunco di judiko. - II. (gram.)}
\l{\cel{3}Asembluro ek objekto, verbo ed atributo. - III. (\sou{geometrio)}.}
\l{\cel{3}Enunco di teoremo demonstrenda. ( IV. (teol.)\cel{2}Afirmo dogma-}
\l{\cel{3}tala. - DEFIRS.}
\l{}
\l{\sou{propra}. Qua proprietesas da persono, da kozo, exkluze di om-}
\l{\cel{3}na altra. - EFIS.}
\l{}
\l{propretoro.\cel{1}\sou{(en Roma, antique}). Lietnanto (remplasanto) di la}
\l{\cel{3}pretoro. - DEFIRS.}
\l{}
\l{proprietar. (trans.)Posedar (ulo) kom sua, kom propra, komplete}
\l{\cel{3}e legitime. - EFIS}
}
